\scanpage{91}
\mspage{78}

il faut donc prouver
\[
\mathrm{TF}'\overset{\mathrm{déf}}{=}C'_*\subset W'.
\]
Cela va résulter du

\uline{Lemme} 4.5.3. On a $\mathrm{TF}'=\unsure{u|v}^{\unsure{*|\deleted{*}}{}^{-1}}(\mathrm{TF})$

C'est formel par définition de $\mathrm{TF}'$ comme
\[
\mathrm{TF}'=\bigl(u_!(\deleted{\unsure{\mathrm T}}C)\bigr)_*,
\]
de $\mathrm{TF}$ comme
\[
\mathrm{TF}=C_*,
\]
et de l'adjonction de $u_!,u^*$. Comme $\mathrm{TF}\subset W$, on en conclut
\[
\mathrm{TF}'=u^{*-1}(\mathrm{TF})\subset u^{*-1}(W)=W',
\]
\unsure{qed.}

Donc on a prouvé ceci :

\uline{Proposition} 4.6. Considérons un couple de foncteurs adjoints
\diagram{\begin{tikzcd}[ampersand replacement=\&,column sep=huge] M \arrow[r,shift left,"u_!"] \& M' \arrow[l,shift left,"u^*"]
\end{tikzcd}\,,}
où $M$ est muni d'une structure de cofibration \uline{spéciale} $(W,C)$. On suppose de plus \unsure{\deleted{que}|{}}les conditions suivantes satisfaites
\[
\left\{\begin{aligned} 1^\circ)&\ M'\ \text{stable par petites }\varinjlim,\ \text{ses éléments sont}\\ &\qquad\text{accessibles}\\ 2^\circ)&\ u^*\ \text{commute aux petites }\varinjlim\\ 3^\circ)&\ u^*(u_!(C))\subset C.
\end{aligned}\right.\tag*{(4.6.1)}
\]
% archive sheet number: 90
